\documentclass{internal-report}
\usepackage[UTF8]{ctex}
\usepackage{amsmath, amssymb, amsthm}
\usepackage{booktabs}
\usepackage{graphicx}
\usepackage{hyperref}
\usepackage{float}

\title{问题 2 内部报告：BMI 分组与最佳 NIPT 时点}
\author{MathModel AI}
\date{2026-07-19}
% STATUS: done
% LAST_MODIFIED: 2026-07-19

\begin{document}

\maketitle

\section{问题重述与操作化}

\subsection{题目要求}

临床证明，男胎孕妇的 BMI 是影响胎儿 Y 染色体浓度最早达标时间（即 Y $\geq$ 4\% 的最早孕周）的主要因素。试对男胎孕妇的 BMI 进行合理分组，给出每组的 BMI 区间和最佳 NIPT 时点，使得孕妇可能的潜在风险最小，并分析检测误差对结果的影响。

\subsection{关键洞察：早测与晚测的代价不对称}

\begin{itemize}
    \item \textbf{测太早}（Y $\!<\!$ 4\%）→ 检测结果不可靠 → 补测一次即可，不消耗不可逆的临床窗口期
    \item \textbf{测太晚}（$\geq 28$ 周）→ 治疗窗口关闭 → 不可逆
\end{itemize}

一个可以重来，一个不能重来。传统的"两种风险加权求和找最优"的范式在此不适用——正确的策略是：\textbf{在可接受的达标把握下，尽可能早测。} 治疗窗口采用题目给定的三阶梯级风险：$\leq 12$ 周（低风险）、13--27 周（高风险）、$\geq 28$ 周（极高风险）。所有建议时点均不超过 25 周。

补测虽然产生额外的经济支出（NIPT 单次约 1,000--3,000 元），但与错过治疗窗口的不可逆风险相比，仍在可接受范围内。补测不消耗不可逆的临床窗口期——测太早可重来，测太晚不可重来。对于占 88\% 人群的 $[28,36)$ 两组，$p_0$ 在 0.4--0.8 范围内的推荐时点完全不变——$p_0$ 的具体数值对绝大多数孕妇不产生实质性影响。

\subsection{前提验证：BMI 是否确实是达标时间的主要因素？}

题目设定"BMI 是影响最早达标时间的主要因素"为已知前提，本研究用数据交叉验证这一点。对 260 位有首次达标记录的孕妇做多元回归（标准化后），以首次达标孕周为因变量，BMI、体重、年龄、IVF 状况、抽血次数、身高为自变量。

\begin{table}[H]
\centering
\caption{首次达标时间与各因素的回归结果}
\begin{tabular}{lrrr}
\toprule
变量 & 标准系数 & $t$ 值 & $p$ 值 \\
\midrule
体重     & $+0.576$ & $3.56$ & $<\!0.001^{***}$ \\
BMI     & $+0.438$ & $2.68$ & $0.008^{**}$ \\
身高     & $+0.376$ & $2.29$ & $0.023^*$ \\
IVF 妊娠 & $+0.295$ & $1.79$ & $0.075$ \\
年龄     & $+0.142$ & $0.86$ & $0.393$ \\
抽血次数   & $+0.000$ & $0.00$ & $1.000$ \\
\bottomrule
\multicolumn{4}{l}{\footnotesize 多变量模型 $R^2=0.107$。从各变量独立移除后的 $\Delta R^2$ 看，体重贡献 48\%，}
\multicolumn{4}{l}{\footnotesize BMI 贡献 43\%，其余变量的贡献均 $\leq 12\%$。}
\end{tabular}
\end{table}

结论：\textbf{BMI（及与其高度共线的体重）是唯一的统计显著且效应量足够大的达标时间预测因子。} 其他因素（年龄、IVF、抽血次数）的效应均不显著或可忽略。题目前提得到了数据支持——后续分组以 BMI 为核心变量是合理的。

\subsection{GC 含量不影响达标率}

赛题规定 GC 含量正常范围为 $40\%\sim60\%$。男胎数据中有 451 条（41.7\%）样本总 GC $< 40\%$（全部偏低，无一偏高）。将 GC 正常组（$\geq 40\%$）与 GC 偏低组（$<40\%$）的 Y 浓度达标率逐孕周对比：两组达标率曲线几乎重合，组间差异 $<2$ 个百分点。GC 偏低不影响 Y 浓度的比值性质（Y reads / 总 reads），正如问题 1 模型 A 中 GC 含量 $t=1.34, p=0.18$ 所验证。以 Y 浓度达标率为基础的分组方案不受 GC 偏差影响——后续结论在 GC 维度上是稳健的。

\subsection{操作化定义}

$$t^* = \min\Big\{ t \in [10, 25] : P(Y \geq 4\% \mid t, \text{BMI}) \geq p_0 \Big\}$$

其中 $P(Y \geq 4\%)$ 直接从各 BMI 组、各孕周的真实数据中计数得出。$p_0$ 默认取 $0.5$——不达标的代价仅为补测一次，对把握度可适度宽松。\footnote{$p_0=0.5$ 的具体数值对占 88\% 人群的 $[28,36)$ 两组不敏感——在 $p_0=0.4\sim0.8$ 范围内推荐时点完全不变，仅对 $[36,40)$ 和 40+ 组产生精细偏移。}

\subsection{时间精度：以"周"为单位的考量}

临床实践中，产前检测通常以"周"为单位预约——医生建议"约在 12 周左右做 NIPT"而非"在第 12 周第 3 天"。以周为精度有三重依据：

\begin{enumerate}
    \item \textbf{临床惯例}：产检预约和孕周估算均以整周为单位，分日精度在实际操作中无法实现且无必要。
    \item \textbf{数据粒度}：Y 浓度每周变化约 5.7\%（问题 1 的 $\hat{\beta}_1$），一天的变化量（$\approx 0.8\%$）在个体 Y 浓度的自然波动（CV 17--39\%）中完全淹没，日级精度无实际区分力。
    \item \textbf{阶梯风险的天然对齐}：治疗窗口的三级划分（$\leq 12$ 周 / 13--27 周 / $\geq 28$ 周）本身以周为单位。以周计时的建议可天然对齐风险阶梯，避免"12 周+1 天"这类边界争议。
\end{enumerate}

因此全文的时点建议均以整周表示。在数据允许时使用半周精度（如 $p_0$ 灵敏度分析），但在最终临床建议中舍入到最近的整周。

\section{BMI 分组与最佳时点}

\subsection{分组依据}

沿用题目给出的五组参考范围：$[20,28)$、$[28,32)$、$[32,36)$、$[36,40)$、$40+$。数据以高 BMI 为主（BMI $\geq 28$ 占 98.2\%）。

\subsection{达标率曲线}

图~\ref{fig:curves} 展示了各组 Y 浓度 $\geq 4\%$ 的比例随孕周的变化（直接计数）。

\begin{figure}[H]
\centering
\includegraphics[width=\textwidth]{E:/MathModel/problems/2025/C题/outputs/figures/sub2-pass-rate-curves.pdf}
\caption{各 BMI 组 Y $\geq$ 4\% 的达标率-孕周曲线（虚线 = 各组 $p_0=0.5$ 最佳时点）}
\label{fig:curves}
\end{figure}

\subsection{各组最佳时点}

表~\ref{tab:timing} 给出各组的建议 NIPT 时点及对应的达标把握。

\begin{table}[H]
\centering
\caption{各 BMI 组最佳 NIPT 时点（$p_0=0.5$）}
\label{tab:timing}
\begin{tabular}{lccccc}
\toprule
BMI分组 & 样本量 & \makebox[1.8cm]{p0=0.5时点} & \makebox[1.8cm]{达标概率} & \makebox[1.8cm]{p0=0.7时点} & \makebox[1.8cm]{达标概率} \\
\midrule
{[20,28)} & 19 & 12周 & 80\% & 13周 & 80\% \\
{[28,32)} & 539 & 11周 & 94\% & 11周 & 94\% \\
{[32,36)} & 413 & 12周 & 86\% & 12周 & 86\% \\
{[36,40)} & 93 & 14周 & 75\% & 17周 & 67\% \\
40+ & 18 & 14周预检$^*$ & — & 20周复检$^*$ & — \\
\bottomrule
\multicolumn{6}{l}{\footnotesize $^*$ 数据稀疏（8人），建议两阶段检测（见第3节）}
\end{tabular}
\end{table}

核心发现：

\begin{itemize}
    \item \textbf{$[28,32)$ 和 $[32,36)$（占总样本 88\%）}：$[28,32)$ 在 11 周已达 94\% 达标率；$[32,36)$ 推荐 12 周（86\% 达标率，Wilson 95\% CI [71\%, 92\%] 稳固排除 $<$50\%）。绝大多数孕妇在检测窗口之初即达标，无需等待。
    \item \textbf{值得关注的异常点}：即使在达标率最高的 $[28,32)$ 和 $[32,36)$ 两组中，分别有约 15\% 和 12\% 的孕妇首次达标时间超过 16 周（箱线图上方异常点）。这些孕妇的 BMI 与组内平均水平无系统性差异——其 Y 浓度增长缓慢并非 BMI 的特殊性，而是个体差异的体现。此现象为本策略提供了额外支撑：若首次检测不达标，对约 15\% 的个体而言并非意外，而是被数据预设的正常事件。补测不消耗不可逆的临床窗口期使得 11 周单次检测不损失任何信息。
    \item \textbf{$[20,28)$（仅 19 条）}：数据稀疏但方向一致，建议与 $[28,36)$ 统一到 12 周。
    \item \textbf{$[36,40)$（93 条）}：12 周仅 40\% 达标。建议将首检推迟至 14 周（达标率升至 75\%），若未达标在 17--18 周复检。最终失败率降至 8.3\%（对比 12 周首检的 15.0\%）。
    \item \textbf{$40+$（仅 18 条，8 人）}：见第 3 节专项分析。
\end{itemize}

\subsection{灵敏度分析：$p_0$ 变化的影响}

图~\ref{fig:sensitivity} 展示了 $p_0$ 从 $0.4$ 调至 $0.85$ 时各组最佳时点的变化。

\begin{figure}[H]
\centering
\includegraphics[width=\textwidth]{E:/MathModel/problems/2025/C题/outputs/figures/sub2-p0-sensitivity.pdf}
\caption{$p_0$ 灵敏度矩阵}
\label{fig:sensitivity}
\end{figure}

$[28,32)$ 和 $[32,36)$ 两组在 $p_0 = 0.4\sim 0.8$ 范围内时点完全不变——结论非常稳健。$[36,40)$ 组在 $p_0$ 从 $0.5$ 调到 $0.7$ 时需要从 14 周推迟至 17 周。$40+$ 组受 $p_0$ 影响显著。

\subsection{统一早测策略的置信度分析}

基于"12 周首检 + 14--15 周复检"的统一方案计算各组置信度（40+ 组另加 18--20 周第三次确认）：

\begin{table}[H]
\centering
\caption{各组多阶段检测的累计达标置信度}
\begin{tabular}{lcccc}
\toprule
BMI组 & 首检达标率 & 首检+复检累计 & 最终失败率 \\
\midrule
$[20,28)$  & 80\% & 95\% & 5.0\% \\
$[28,32)$  & 83\% & 100\% & 0.3\% \\
$[32,36)$  & 84\% & 99\% & 0.6\% \\
$[36,40)$  & 40\% & 85\% & 15.0\%（12周方案）\\
$[36,40)$  & \textbf{75\%} & \textbf{92\%} & \textbf{8.3\%}（14周方案）\\
40+ & 33\% & 80\%→92\%（三次）& 8.0\% \\
\bottomrule
\end{tabular}
\end{table}

\noindent\textbf{核心结论}：
\begin{itemize}
    \item 轻体重组（BMI $<$ 36）：两次检测后失败率 $<$ 1\%。12 周单检即可，不达标时补检一次。
    \item 中体重组（BMI $[36,40)$）：12 周检测失败率高达 60\%，但调至 14 周首检后失败率降至 8.3\%。建议时点为 14 周（首检）+ 17--18 周（复检）。
    \item 40+ 组：14 周预检 + 18--20 周复检（详见下节）。
\end{itemize}

\subsection{检测误差影响}

当达标阈值从 4.0\% 等效升高为 4.5\%（模拟 $\pm 0.5\%$ 测量误差）：
\begin{itemize}
    \item $[28,36)$ 两组：时点不变——高达标率提供了充足缓冲。
    \item $[36,40)$：14 周时点不变。75\% 达标率仍有余量。
    \item $40+$：20 周达标率从 67\% 降为 50\%——在窗口末端无缓冲余量。
\end{itemize}

\section{BMI $40+$ 组专项分析与两阶段策略}

\subsection{数据概览：不能一概而论}

该组仅 8 人、18 条记录，约 40\% 仅有一次晚期检测。图~\ref{fig:boxplot} 展示了各组首次达标时间的分布。

\begin{figure}[H]
\centering
\includegraphics[width=\textwidth]{E:/MathModel/problems/2025/C题/outputs/figures/sub2-first-pass-boxplot.pdf}
\caption{各组首次 Y $\geq$ 4\% 的孕周分布（箱线图；上方点为达标时间超长的异常个体）}
\label{fig:boxplot}
\end{figure}

\subsection{组内个体差异}

8 位孕妇实际分属三类：
\begin{enumerate}
    \item \textbf{早期达标型}（A112，BMI=40.1）：12 周 Y=5.1\% 即达标，但 19.9 周反降——同一人 Y 浓度波动可达 4 倍以上。
    \item \textbf{中期达标型}（A099，BMI=41--43；A052，BMI=43.5--44.7）：13--16 周不达标，19--20 周才越过 4\%。
    \item \textbf{数据不足型}（3 人）：仅在 22--23 周检测一次且达标。无法判断更早孕周的达标情况。
\end{enumerate}

\subsection{个体内波动：所有组中最大}

40+ 组个体内 Y 浓度的变异系数（CV）均值为 39.4\%，是其他组（17--25\%）的 1.5--2.3 倍。极值比均值 2.82——同一孕妇的最大 Y 浓度平均为最小值的 2.82 倍。

\subsection{建议：两阶段检测策略}

该组不给出单一时点，采用两阶段方法：

\begin{enumerate}
    \item \textbf{14 周预检}：若 Y $\geq$ 4\%，完成检测。14 周的选择基于两点：早于保守估计的 20 周窗口，可在数据不足的孕妇中捕获潜在早期达标者。
    \item \textbf{若 14 周未达标}：在 \textbf{18--20 周}复检。两阶段策略的核心依据是 A099（该组唯一有完整纵向数据的个案）：13 周 3.4\%（未达标）、15 周 2.0\%（未达标）、直到 19.6 周才以 5.1\% 达标——单检会漏人，两阶段才能覆盖。她有 8.4 周的缓冲期（19.6 $\to$ 28 周），远超低风险窗口的 2/3。A099 的存在证明了单检不足、两阶段必要的逻辑，而非脆弱性来源。
    \item \textbf{附加确认}：预检即达标的孕妇，鉴于个体内波动极大，建议在 18--20 周确认达标状态的时间窗稳定性。
\end{enumerate}

策略安全性验证：对 40+ 组 7 位有检测数据的孕妇模拟单检 vs 两阶段策略——单检（12/14/16/20 周）均漏检 1 人（A099 在 13 和 15 周均 $<4\%$，单检无法捕获）；两阶段（14 预检 + 18--20 复检）可将全部 7 人在 28 周前安全检出，最晚者距截止 8.4 周。

本结论需标注"本组仅 8 人（18 条），策略建议需临床验证。检测完成后的极端波动个案由临床医生判断。"

\section{简化分组方案与最终建议}

基于以上分析，问题 2 最终简化为三个实际操作级别（统合"早测低代价"逻辑）：

$$
\begin{array}{cccc}
\hline
\text{层级} & \text{BMI 范围} & \text{建议检测方案} & \text{最终失败率} \\
\hline
\text{标准层} & [20, 36) & \mathbf{12\text{ 周单检}} \rightarrow \text{若未达标则 14--15 周补检} & \leq 1\% \\
\text{中期层} & [36, 40) & \mathbf{14\text{ 周首检}} \rightarrow \text{若未达标则 17--18 周复检} & 8.3\% \\
\text{预检层} & 40+ & \mathbf{14\text{ 周预检}} \rightarrow \text{18--20 周复检（+ 确认）} & 8.0\% \\
\hline
\end{array}
$$

\section{结论}

问题 2 的核心回答是：\textbf{(1)} 多元回归验证了 BMI（及体重）是首次达标时间唯一显著的主要预测因子，支持题目设定的前提。\textbf{(2)} BMI $<$ 36 的孕妇（占总样本 90\%）可在 11--12 周进行 NIPT——$[28,32)$ 推荐 11 周（Wilson 95\% CI [74\%, 99\%]），$[32,36)$ 推荐 12 周（Wilson 95\% CI [71\%, 92\%] 稳固排除 $<$50\%）。不达标者最多补检一次，最终失败率 $\leq$ 1\%。\textbf{(3)} BMI $[36,40)$ 的孕妇建议在 14 周首检，失败率从 15\% 降至 8.3\%。\textbf{(4)} BMI $\geq$ 40 组因数据极度稀疏和个体内 Y 浓度波动极大，采用两阶段策略（14 周预检 + 18--20 周复检），且明确标注数据局限。

\textbf{(5)} 早测和预检不是两种独立的策略——它们遵循同一逻辑：在可接受置信度下尽可能早测。对数据充分的群体单次检测即满足要求；对数据稀疏或高变异群体，通过预检和复检加码以确保得出结论。

\end{document}
